% =============================================================================
% Revision History
% =============================================================================

\section*{Revision History}
\addcontentsline{toc}{section}{Revision History}

\begin{tabularx}{\textwidth}{@{}llXX@{}}
\toprule
\textbf{Version} & \textbf{Date} & \textbf{Author} & \textbf{Summary}\\
\midrule
\midrule
0.2.4 & 2026-07-30 & Geovana Neves &
When a declared uncertainty applies. A PATCH revision by the rule
chapter~9 states: a clarification, adding no requirement and withdrawing
none. Section~\ref{sec:itceq} now states the moment an
\code{[uncertainties]} declaration takes effect, which it had left
unspecified: before the first equation for a name the file does not
produce, and immediately after the writing equation for a name it does,
keyed on what the FILE produces and never on what the VarFrame carried on
the way in. The implementation had adopted the other rule, so a target the
frame already carried was assigned before the loop and then overwritten by
its own equation's propagation: one file gave two answers, and in the
worse branch it shipped \code{u(y) = 2.0} where the file declares
\code{5.0}. A wrong number rather than a missing one, chosen by nothing the
caller can see (\code{R4-ITA-003}, \code{ITC-20260730-0105}).
\\[4pt]
0.2.3 & 2026-07-30 & Geovana Neves &
The sixth CHK-1 refusal, and the build. A PATCH revision by the rule
chapter~9 states: a clarification, adding no requirement and withdrawing
none. REQ-01 gains the refusal of a CSV row WIDER than its header
(\code{DataError}, \code{R3-ITA-009}) and states the asymmetry out loud,
that a row NARROWER than its header is still NaN-filled, since a reader
who meets one half of it infers the other wrongly. The asymmetry is
deliberate: a row with more fields than the header is a malformed file,
either a truncated header or a delimiter inside an unquoted field, while a
row with fewer is ordinary missing data. REQ-01 also states the
empty-cell rule the asymmetry leans on, which the document had nowhere.
This refusal shipped in v0.2.0 marked BREAKING in \code{CHANGELOG.md} and
its own error message cited REQ-01, while REQ-01 did not describe it; the
0.2.2 revision below closed five such gaps and missed this one, which the
independent review REV-004 found (\code{ITC-20260729-2255}).

REQ-95 gains the specification build. Two defects had reached \code{main}
in a document that no mechanism compiled: chapter~7 carried a blank line
after every content line, and one unescaped underscore inside \code{code}
put chapter~8 into math mode. Seven \code{pdflatex} errors, and the only
thing that ever found either was a person running the build by hand, with
\code{main.log} dated 2026-07-23 while 0.2.1 and 0.2.2 both shipped after
it (\code{ITC-20260729-2300}, \code{ITC-20260729-2350},
\code{ITC-20260730-0010}). The chapter~7 reflow changed no content: 692
source lines became 346 with the 307 non-blank lines byte-identical and in
the same order.
\\[4pt]
0.2.2 & 2026-07-29 & Geovana Neves &
CHK-1 refusal amendment (\code{ITC-20260729-1450}), which blocked the
v0.2.0 tag. The revision below predicted this class in its own words and
closed four instances of it; the CHK-1 release checkpoint then added
five more public refusals whose requirement text did not describe them,
and this revision closes those. The code shipped correct in every case:
what was wrong was the text, so no requirement scope is added or
withdrawn. REQ-01 states that one name may not be both a dimension and a
variable, so a column named \code{datapoint} is refused
(\code{R3-ITA-008}). REQ-24 and REQ-107 state the \code{itc.concat}
\code{AxesError} on inputs whose vector groups sit in different axis
systems (\code{CHK1-003}). REQ-39 requires a finite declared
uncertainty magnitude, cites GUM clause~5.1.2 by its two operative
sentences, and states the limit of the check (\code{R3-ITA-007},
OQ-40). REQ-44 refuses a built-in constant shadowing a measured channel
and REQ-45 moves from two refusals to three (DD-42, \code{CHK1-001},
OQ-41).

Two enumerations were false in normative text and are corrected
differently, because only one of them should exist at all. REQ-92
claimed \code{0.2.0} ships three breaking changes and it ships more; its
enumeration is REMOVED rather than recounted, and no replacement count
is put in its place, since \code{CHANGELOG.md} is where the same
requirement already mandates the list, a duplicate inventory goes stale
independently of its original, and a bare count is that same artifact
with its evidence removed. REQ-98 carried its provisional family at
three different counts across four places while
Table~\ref{tab:unc-ops} still claimed \code{smooth} and \code{diff}
propagate; the list is now given ONCE, at five, adding \code{fill} with
\code{method="polyfit"}, and REQ-26 gains the matching per-method split.
The ground recorded for \code{fill(method="polyfit")} is that the fill
path emits no weights for the propagation to consume. Explicitly not
that its weights are data dependent, which was measured false and would
have drawn \code{interpolate} into the family with it; and explicitly
not that no polynomial weight rule exists, since one ships and
\code{interpolate} propagates through it. Whether that matrix is
admissible over a support chosen per gap is OQ-42, opened here, and it
may return the count to four.

This is a document-version BUMP rather than a same-baseline refinement,
for the reason the 0.2.1 entry established and one more. Nine \stable{}
requirements change text, the six listed above for the refusals and
REQ-92, REQ-98 and REQ-26 for the two enumerations, which is well past
the threshold that entry set; and a reader holding v0.2.0 needs to be
able to name which document version first described the refusals that
release actually shipped. The argument does not rest on the count being
exact, which is why it is given with its list beside it.\\

0.2.1 & 2026-07-28 & Geovana Neves &
Independent-review amendment (REV-001, reproduced as \code{ITA-0}).
REQ-103 is restated as a SEMANTIC GUARANTEE rather than an enumeration
of named fields: two VarFrames in the same semantic state have the same
hash, and the field list follows from that definition instead of
constituting it. The scope gains dimension and variable metadata (unit,
description, long name) and names the axis registry, closing a drift
that ran in BOTH directions inside one requirement, since the registry
was already hashed by the code while the reqbox omitted it and units
were outside the hash while \code{rotate} read them (\code{ITACA-003}:
a \code{deg} frame and a \code{rad} frame hashed identically and
produced \code{FZ} of $-1.0$ against $-0.894$). Same semantic state is
defined at the floating-point boundary in DD-40. REQ-107 moves from
draft to stable in the same revision, so a stable guarantee stops
taking an input whose token format a draft requirement defines, and it
is amended to record that \code{db.rotate} UPDATES a group's axis.
REQ-36's guard is stated as per-subtree rather than per-expression
(\code{ITACA-022}), and REQ-45 records that \code{validate} refuses a
\code{[constants]} name the frame carries as a measured variable
(\code{ITACA-002}, DD-39).

Four further stable requirements are amended, each because this lane
added a public refusal or a state rewrite to a surface whose text did
not describe it. REQ-40 gains the positive-semidefinite condition on
the assembled matrix (\code{ITACA-001}). REQ-38 gains the correlation
write-back from the transformed covariance and its three refusal
conditions (\code{ITACA-025}). REQ-44 states that expression calls
take positional arguments only (\code{ITACA-023}). REQ-71 and P-02
declare the two in-memory export bridges deliberately lossy
(\code{ITACA-005}). Chapter~8's false PROV-N claim
(\code{ITACA-012}) and its overstated docstring-enforcement claim
(\code{ITACA-016}) are corrected in the same revision.

This is a document-version BUMP rather than a same-baseline refinement,
unlike the amendments folded into 0.2.0 above, for one reason: it is
the first amendment that invalidates artifacts a published release
wrote. An \code{.itc} archive from v0.1.0 carrying units no longer
verifies, and a reader needs to be able to name which document version
changed the rule.\\

0.2.0 & 2026-07-23 & Geovana Neves &
M1 re-baseline (ultraplan Batch A): the Chapter 10 M1 deliverables
split into a v0.2.0 computation table and a v0.2.1 presentation and
builtins table, with the stretch scope allowed to ship inside v0.2.0
if green in the same window. Library-infrastructure requirements
added as draft from the 2026-07-23 library-review triage: options
registry with exact keys (REQ-104, DD-24), typed no-default sentinel
(REQ-105), accessor registration (REQ-106). Design-decision boxes
DD-23 to DD-25 enter Chapter~5 with their long-form entries in the
decision log; DD-25 records the dev-only \code{uncertainties} test
oracle (no requirement). REQ-101 was validated as written at the same
checkpoint; its status moves to stable through the normal process
once implemented and tested during M1. At the M1 Phase B1
checkpoint, REQ-31 and the normative UncFrame-semantics table gained
the \code{fitmodel} raise-until-frozen note and row (OQ-24), and
OQ-24 and OQ-25 were opened. Three author-seat calls at the same
checkpoint: \code{fitvalue} defers with \code{fitmodel} and raises on
uncertainty (OQ-24); the too-few-points-for-degree invariant is
unified into a shared \code{FitDegreeError} leaf across \code{diff},
\code{smooth}, \code{interpolate}, and \code{fitmodel} (REQ-30, the
former \code{DiffWindowError}); and \code{fill}'s positional
\code{method} is deprecated toward keyword-only (REQ-26). These are
same-baseline refinements recorded here rather than a further version
bump. M1 Phase B2 opened the axes work: the built-in wind and
stability direction-cosine convention (AIAA R-004A, Etkin) was
SME-accepted; scipy joined the dev-only test oracles (DD-26);
condition-dependent angle inputs are read in the unit of their
Dimension or Variable; the axis registry and vector-group declarations
join the VarFrame state hash (REQ-103); and REQ-107 (per-group source
frames and multi-frame data) was added as a draft requirement. With the
axes machinery implemented and tested (rotation with per-grid-point
condition-dependent frames and chain-rule angle uncertainty, moment
transfer), REQ-101 was promoted from draft to stable. The angle
independence model of the rotation propagation was recorded as OQ-26,
and DD-27 records the sanctioned accessor per-instance cache write.
OQ-25 (origin-tag reduction across collapsing and reshaping
operations) was resolved and its four rules folded into Section 4.3.
REQ-108 (published documentation site on GitHub Pages, mirroring
pyflightstream, DD-23) was added as a draft requirement. M1 Phase B3a
implemented the pipeline requirements and corrected the specification
where it was wrong: Section~\ref{sec:itc-pipe} now specifies the
\code{.itc\_pipe} serialization as JSON rather than TOML, with the reasons recorded in
DD-28 (no standard-library TOML writer in any Python version, no TOML
null type against a meaningful \code{compute(fill=None)}, and nested
replay arguments); every content item the section required is kept,
including the attached comment (REQ-19) and the content hash, which is
now reverified on load. REQ-53 gained the normative extraction contract
(frame construction is skipped as input preparation, any other
step-less operation raises, and a range yielding no step raises rather
than producing a silent no-op), and REQ-54 gained the replay contract
(the recorded call and comment reproduce the state hash when replayed
onto the frame the range was lifted from, while a different frame
reproduces the processing rather than the hash, and an incompatible
target raises with the underlying object and fix preserved). The History-entry table gained the \code{step} field, replay
metadata excluded from the state hash and persisted in the \code{.itc}
archive at schema \code{itaca-itc/2}. Section~\ref{sec:itc-archive} was
amended to match: it now names both schema versions, the \code{step}
member of \code{history.json}, and the \code{steps\_hash} digest in
\code{metadata.json}, with the scope boundary against the REQ-103 state
hash stated there and the reasoning recorded in DD-30. The scope of
REQ-103 is unchanged. M1 Phase B3b amended \stable{} REQ-83: the
language baseline rises from \code{>=3.10,<3.14} to
\code{>=3.11,<3.14}, so \code{tomllib} is in the standard library and
the \code{.itceq} parser of REQ-48 reads the TOML structure of
Section~\ref{sec:itceq} with no dependency, no vendored reader, and no
hand-written parser. This resolves OQ-28, which had blocked the phase
since 2026-07-23, by taking none of the three options it offered;
DD-32 records why, and why it leaves DD-28 and the JSON
\code{.itc\_pipe} encoding standing. REQ-83 also gained the normative
statement that the baseline may not fall below a standard-library
module the implementation depends on, and the note that its four
restatements (the PyPI classifiers, the \code{ruff} target version, the
CI matrix floor, and the imports themselves) are pinned to one
declaration by \code{tests/test\_python\_floor.py}. Phase B3b
implemented REQ-45 to REQ-48 in the new \code{pproc/} package, and
amended two further \stable{} requirements at the same checkpoint.
REQ-82 now states the NumPy-only scope as every package except
\code{io/} and \code{utils/}, rather than naming the three that existed
at the 0.1.0 baseline: enforcement already had that scope, and stated
as an exception list a package added later is restricted by default
instead of inheriting nothing (DD-33). REQ-46's factory is renamed from
\code{itc.pproc} to \code{itc.processor}, resolving OQ-29: the former
name shadowed the \code{pproc} package on that attribute, so the
\code{pproc.statistics} and \code{pproc.compare} forms of REQ-49 and
REQ-50 could never have resolved. Renaming the factory rather than the
package leaves every module path and both requirement texts true as
written (DD-34). The Chapter~\ref{ch:contributing} worked example was
amended with it, since it used the same token for a third thing, the
processor instance. Three author-seat calls closed the phase, each
amending a \stable{} requirement or its section rules. \stable{} REQ-47
now recognizes a reapplication from two pieces of evidence together,
the produced variables being present \emph{and} the History recording
the processor, because matching names alone refused a first
application over inconveniently named channels and taught the caller to
pass \code{force=True} by reflex (DD-35). REQ-47 also admits
\code{[meta] idempotent} as a declaration, so a file-defined processor
no longer has to be subclassed in Python to say that reapplying it is
meaningful (DD-36). Section~\ref{sec:itceq} gained three section rules
the parser enforces: a name declared in \code{[constants]} may not also
be an equation or correction target, because a constant is substituted
into every read and such an equation would run unreachable (DD-37); a forward
reference in file order is refused when the file is read; and
\code{[corrections]} may replace a variable the VarFrame supplies, not
only one \code{[equations]} produced, which is the ordinary wind tunnel
case of correcting a measured channel in place.\\
\addlinespace
0.1.1 & 2026-07-21 & Geovana Neves &
REQ-98 and REQ-99 promoted from draft to stable after the M0
implementation checkpoint: propagation semantics validated for every
M0 operation, with the smooth and diff row of REQ-98 remaining
provisional pending OQ-18. Declared correlation applies to both
uncertainty components (Section 4.2, OQ-23). Folded the OQ-19 to
OQ-22 resolutions into the text: dict-mode \code{"*"} swept sentinel
(REQ-01, REQ-03), Provenance \code{source\_coords} backing the
manifest (Section 4.4.1), string columns as dimension coordinates
only (REQ-01), and the draft-guard scope (REQ-11).\\
\addlinespace
0.1.0 & 2026-07-21 & Geovana Neves &
First workspace-tracked version. Consolidates the founding v0.1.0b draft and
incorporates the 2026-07-21 specification review: uncertainty semantics for
every operation (REQ-98), two-component uncertainty (REQ-99), moment
reference-point transfer (REQ-100), condition-dependent axes (REQ-101),
array immutability (REQ-102), state-hash reproducibility contract (REQ-103),
solver-automation boundary (NREQ-10), design decisions DD-14 to DD-22,
positioning among existing tools, license declaration, and an incremental
release roadmap. Requirement evolution is monitored from this baseline
onward: every change is recorded in this table and in
Chapter~\ref{ch:changelog}, and deprecated requirements are kept in place
and tagged \deprecated{}.\\
\bottomrule
\end{tabularx}

\vspace{1em}
\begin{noteinfo}[About this document]
This Software Requirements Specification is part of the \itc{} project's
public design record. It is being prepared in parallel with the implementation,
which is itself developed with assistance from large language model tooling
(notably Claude Code). Every implementation derived from these requirements
goes through human double-check by the author before being merged. The SRS is
the authoritative reference; if a discrepancy with the code is detected, the
code is corrected, not the SRS.
\end{noteinfo}

\clearpage
